\iffull
Consumer headphones and smartwatches already carry the sensors needed for hands-free laptop control, but turning noisy wrist signals into dependable commands is harder than classifying isolated gestures. We present a volume-control system in which the user looks at an on-screen target, pinches to grab it, turns the wrist to set the level and opens the pinch to let go. Head orientation comes from Sony headphones, photoplethysmography (PPG) and inertial signals come from a Samsung Galaxy Watch, and all inference runs on the laptop. Separate pinch-start and pinch-release detectors, explicit cancellation paths and Off/Monitor/Live modes stop a missed or spurious detection from becoming an unwanted command. We evaluate detection over whole recordings rather than pre-cut windows, counting misses, late and repeated events, and false activations per hour of ordinary activity. On eight held-out recordings from one participant, the pinch-in detector caught 47 of 57 annotated starts with a median latency of 442\,ms and produced no false activations in 0.48\,h of no-gesture time (95\% upper bound: 6.2 per hour). An earlier camera-labelled model caught 161 of 372 starts across days, with large day-to-day variation. We also find that 37 exported recordings contain only 28 distinct sensor streams, a duplication that can quietly leak data between training and test sets. Because the technique needs no camera view of the hand, it may also suit augmented and virtual reality headsets, where the hands are often outside the cameras' view; we have not yet tested it there. We close with the multi-participant, interaction and headset studies needed before broader claims can be made.
\else
Consumer headphones and smartwatches already carry the sensors needed for hands-free computer control, but turning noisy wrist signals into dependable commands is harder than classifying isolated gestures. We present look--pinch--turn: the user looks at an on-screen target, pinches to grab it, turns the wrist to set a value and opens the pinch to let go. Head orientation comes from Sony headphones and photoplethysmography (PPG) and inertial signals from a Samsung Galaxy Watch, with all inference on the laptop. Separate pinch-start and pinch-release detectors and explicit cancellation paths stop a missed or spurious detection from becoming an unwanted command. Evaluated over whole recordings from one participant, the pinch-in detector caught 47 of 57 annotated starts with a median latency of 442\,ms and no false activations in 0.48\,h of no-gesture time (95\% upper bound: 6.2 per hour), while an earlier camera-labelled model varied widely across days. We also show that duplicate exports can quietly leak data between training and test sets. Because a wrist-sensed pinch needs no camera view of the hand, the technique may also suit AR/VR headsets; we outline the studies needed to test this.
\fi
